%! suppress = EscapeUnderscore
Despite their usufullness, RNNs have a major practical, computational, drawback. Namely, since they are sequential in
nature, they are dealing with a large number of training steps and due to their recurrent nature they are not
parallelizable. This makes them very slow to train, since we have to wait for the previous element to be processed
before we can process the next one, due to the recurrent relation between steps. \v

A solution to this problem, was introduced in 2017 by Vaswani et al.\ in the paper ``Attention Is All You Need''. The
solution is a model called ``transformer''.

\bd[Transformer]
\textbf{Transformer} is a neural network architecture that uses attention mechanism to draw global dependencies
between input and output.
\ed

The transformer architecture is based on the encoder-decoder architecture, which we have already seen in the context
of RNNs. However, in contrast to RNNs, the transformer architecture does not use any recurrent units, but instead it
uses attention mechanism to draw global dependencies between input and output.

\fig{trf01}{0.175}

In what follows we will go step by step through the transformer architecture, and we will dive into the details of
each component.

\section{Input Embedding}

\fig{trf02}{0.3}

The first step of the transformer architecture is the so called ``input embedding''.

\bd[Input Embedding]
\textbf{Input embedding} is a simple embedding layer that takes as input the input sequence, and it outputs a vector
representation for each element of the input.
\ed

Transformers are mainly used for natural language processing (NLP) tasks, hence the input sequence is usually a
sequence of words. The input embedding layer takes as input a sequence of words, and it outputs a vector
representation for each word. For this reason we will spend some time on the concept of word embedding, since it is
very important for NLP tasks. \v

In natural language processing (NLP), word embedding is a term used for the vectorial representation of words for
text analysis, typically in the form of a real valued vector  Word embeddings can be obtained using a set of language
modelling and feature learning techniques where words or phrases from the vocabulary are mapped to vectors of real
numbers. Conceptually it involves a mathematical embedding from a space with many dimensions per word to a continuous
vector space with a much lower dimension. \v

\bd[Word Embedding]
\textbf{Word embedding} is the collective name for a set of language modelling and feature learning techniques in
natural language processing (NLP) where words or phrases from the vocabulary are mapped to vectors of real numbers.
\ed

Let us start with a very simple motivation for why we are interested in a vectorial representations of words. Suppose
we are given an input stream of words like a sentence or a document, and we are interested in learning some function
of it e.g: $\hat{y} = \text{sentiments(words)}$. Say, we employ a machine learning algorithm (some mathematical
model) for learning such a function $\hat{y} = f(x)$. We first need a way of converting the input stream (or each
word in the stream) to a vector. Let's start with two important definitions. \v

\bd[Corpus]
A \textbf{corpus} is a language resource consisting of a large and structured set of texts.
\ed

\bd[Vocabulary]
Given a corpus, a \textbf{vocabulary} $V$ is the set of all unique words in the corpus.
\ed

There are two big classes of methods of word embedding: the ``count based models'' which they use the co-occurrence
counts of words, and the ``prediction based models'' which directly learn word representations. The first class of
count based models were the one used before the discovery of deep learning while prediction based models are the
latest advancements in the area of word embedding that make use of deep learning. We will begin by briefly
introducing the count based models, since they will be useful for later, and then we switch to the main part of
prediction based models.

\subsection{One-Hot Representation}

\bd[One-Hot Representation]
Given a vocabulary $V$ with $n$ terms, the \textbf{one-hot representation} is a vectorial representation of each
term in the vocabulary, where each vector is of dimension $n$, and it has only one element equal to 1, and the rest
equal to 0.
\ed

\be
Here is an example of a one-hot representation of the vocabulary $V = [$ cat, dog, truck $]$:

\fig{onehot}{0.25}
\ee

While one-hot representation is very simple and intuitive, it carries a lot of problems. First of all, usually
vocabularies tend to be very large hence the vector space turns enormous which is computationally very inefficient.
Second, the representation is very sparse, meaning that the vast majority of the entries are zeros. Third, and most
important, the representation does not capture any notion of similarity among the words. In our example, ideally, we
would want the representations of cat and dog to be closer to each other than the representations of cat and truck.
However, with one-hot representation, the Euclidean distance between any two words in the vocabulary is simply
$\sqrt{2}$. Similarly, the cosine similarity between any two words in the vocabulary is $0$.

\subsection{Distributed Representation}

We will now introduce the distributional similarity based representation. The idea behind distributed representation
is expressed in the phrase ``you shall know a word by the company it keeps''.

\bd[Co-Occurrence Matrix]
Given a vocabulary $V$ with $n$ terms, the \textbf{co-occurrence matrix} is an $(n \times n)$ matrix which captures
the number of times a term appears in the context of another term, where context is defined as a window of $k$ words
around the terms. Each row (or column) of the co-occurrence matrix gives a vectorial representation of the
corresponding word.
\ed

\be
As an example, let us build a co-occurrence matrix for a toy corpus with $k = 2$. The corpus consist of the following
sentences:
\bit
\item Human machine interface for computer applications.
\item User opinion of computer system response time.
\item User interface management system.
\item System engineering for improved response time.
\item System engineers optimize for human experience.
\eit

The vocabulary is: $V = [$ Human, Machine, Interface, For, \ldots $]$, and the co-occurrence matrix:

\fig{cooccurance}{0.4}
\ee

Of course, as with one-hot, also this representation carries its own problems. First of all, stop words like (``a'',
``the'', `` for'', \ldots) are very frequent hence these counts will be very high. This however, is easily solvable
since we can either ignore very frequent words or simply set a threshold $t$ (for example $t=100$) and whenever the
count of a word gets more that $t$ we simply stop counting further, and we use the threshold $t$ as the value,
ignoring all the other occurrences. \v

The most important problems of distributed representation though are that, as in hot-one representation, it is very
high dimensional, it is very sparse, and it grows with the size of the vocabulary. All of these problems were solved
by using singular value decomposition as a dimensionality reduction technique. However, we will not see that in more
detail since it is not related to deep learning, and it gets out of topic. \v

For many years the singular value decomposition of distributed representation was the way to go for NLP and word
embedding problems. However, in 2006 and after, after the discovery of deep learning, we switched gears from the
whole idea of counting word occurrences and co-occurrence matrices of count based models to direct, prediction based
models that uses statistics to predict outcomes. In what follows we will introduce the most heavily used group of
techniques of prediction based models called ``Word2vec''.

\subsection{Word2vec}

\bd[Word2vec]
\textbf{Word2vec} is a technique for natural language processing published in 2013, that uses a neural network model to
learn word associations from a large corpus of text and to detect synonymous words or suggest additional words for a
partial sentence.
\ed

As the name implies, Word2vec represents each distinct word with a particular list of numbers called a vector. The
vectors are chosen carefully such that a simple mathematical function (the cosine similarity between the vectors)
indicates the level of semantic similarity between the words represented by those vectors. Embedding vectors created
using the Word2vec algorithm have some advantages compared to earlier algorithms that we described previously. \v

Word2vec is a group of related models that are used to produce word embeddings. These models are shallow, two-layer
neural networks that are trained to reconstruct linguistic contexts of words. Word2vec takes as its input a large
corpus of text and produces a vector space, typically of several hundred dimensions, with each unique word in the
corpus being assigned a corresponding vector in the space. Word vectors are positioned in the vector space such that
words that share common contexts in the corpus are located close to one another in the space. \v

As we mentioned, Word2vec can utilize either of two model architectures to produce a distributed representation of
words:
\bit
\item \textbf{Continuous bag-of-words (CBOW)}: Predicts the current word from a window of surrounding context words.
The order of context words does not influence prediction (bag-of-words assumption).
\item \textbf{Continuous skip-gram}: Uses the current word to predict the surrounding window of context words. It
weighs nearby context words more heavily than more distant context words.
\eit

According to the authors' note CBOW is faster while skip-gram does a better job for infrequent words. In what follows
we will introduce both of them since they are the standard way of dealing with word embeddings today.

\subsubsection{Continuous Bag-Of-Words}

\bd[Continuous Bag-Of-Words (CBOW)]
\textbf{Continuous Bag-Of-Words} (\textbf{CBOW}) is a model that tries to predict a target word based on the context of
surrounding words.
\ed

CBOW  tries to understand the context of the words and takes this as input. It then tries to predict words that are
contextually accurate. In other words it tries to predict the current target word based on the source context words (
surrounding words). Let's see an example.

\be
Consider the simple sentence ``the man sat on a chair''. This can be seen as pairs of (context, word) where if we
consider a context window of size 1 (for simplicity) we have ([the], man), ([man], sat), ([sat], on), ([on], a), ([a
], chair). We have used a list for context just to make obvious that if one picks a window greater than 1 then we
would have more context words, e.g.\ for a window of size 2 we would example like ([the, man], sat), ([man, sat], on
) and so on \ldots Thus the model tries to predict the target word based on the context word.
\ee

The CBOW model solves this problem by using a feedforward neural network. The input to the network will be a one-hot
representation of the context words and the output, will be a probability distribution over all possible words in the
vocabulary (multi-class classification problem).

\be
Let's see this in our simple example with the window of size 1.

\fig{cbow}{0.45}
\vspace{-10pt}

In this particular example the input layer to the network is a one-hot representation of the context word ``sat''
which is:

\bse
\text{``sat''} = [0,1,0,0, \ldots, 0]
\ese
\ee

The CBOW network is shallow since it contains only one single hidden layer. The first set of weights $W^{[1]}$ are of
course of dimensions $(k \times |V|)$ where $k$ is the number of units in the hidden layer and $|V|$ is the size of
the vocabulary which is of course also the size of the one-hot representation vectors. It is common practice to call
$W^{[1]}$ simply $W_{\text{context}}$. Notice that we carry no biases. The pre-activations are:

\bse
\textbf{a}^{[1]} =W_{\text{context}} \cdot \textbf{x}
\ese

and the activation function $g$ is as simple as $g(\textbf{x}) = \textbf{x}$, hence the activation of the hidden
layer $\textbf{h}^{[1]}$ is:

\bse
\textbf{h}^{[1]} =W_{\text{context}} \cdot \textbf{x}
\ese

Notice that the product $W_{\text{context}} \cdot \textbf{x}$, given that $\textbf{x}$ is a one hot-vector, is
simply the $\text{i}^{\text{th}}$ column of $W_{\text{context}}$ since everything else will be multiplied by 0.
So when the $\text{i}^{\text{th}}$ word is present the $\text{i}^{\text{th}}$ element in the one-hot vector
is on and the $\text{i}^{\text{th}}$ column of $W_{\text{context}}$ gets selected. In other words, there is a
one-to-one correspondence between the words and the column of $W_{\text{context}}$. More specifically, we can treat
the $\text{i}^{\text{th}}$ column of $W_{\text{context}}$ as the representation of the
$\text{i}^{\text{th}}$ context.

\be
\vspace{-15pt}
\fig{cbow2}{0.3}
\ee

Moving on, the weights of the output layer $W^{[2]}$ are of course of dimensions $(|V| \times k)$, since we need an
output (aka a probability) for each possible word in the vocabulary (aka a probability distribution). It is common
practice to call $W^{[2]}$ simply $W_{\text{word}}$. The pre-activations now are:
\bse
\textbf{a}^{[2]} =W_{\text{word}} \cdot \textbf{h}^{[1]}
\ese

and for the output function $O$ it makes sense to use a softmax since we are dealing with a multi-class
classification problem, hence:
\bse
\hat{\textbf{y}} = P(\text{word \:} | \text{\: context}) = \frac{e^{W_{\text{word}} \cdot
\textbf{h}^{[1]}}}{\sum_{\text{words}} e^{W_{\text{word}} \cdot \textbf{h}^{[1]}}}
\ese

As we can see, since:
\bse
W_{\text{word}} \cdot \textbf{h}^{[1]} = W_{\text{word}} \cdot W_{\text{context}} \cdot \textbf{x}
\ese

\v

and since from the product $W_{\text{context}} \cdot \textbf{x} \:$ survives only the column of
$W_{\text{context}}$ associated with the context, it turns out that $P (\text{word \:} | \text{\: context})$ is
proportional to the dot product between of $W_{\text{context}}$ and $W_{\text{word}}$. In other words the
probability for each word depends on the $\text{i}^{\text{th}}$ column of $W_{\text{word}}$. We thus treat the
$\text{i}^{\text{th}}$ column of $W_{\text{word}}$ as the representation of the $\text{i}^{\text{th}}$ word.
\v

Now the forward propagation is done and the network is built. In order to train it we move to backpropagation. First
of all, we need to calculate the loss function. Since we are dealing with a multi-class classification problem we
will use the cross entropy loss function. We have already calculated the cross-entropy loss function for logistic
regression. The calculations for a softmax regression are very similar, so we will skip them. Just for completeness
we will use the principle of maximum likelihood to obtain the cross entropy loss function for softmax which simply
is:
\bse
\mathcal{L}(\theta) = - \ln P(\text{word \:} | \text{\: context}) = - \ln \frac{e^{W_{\text{word}}
\cdot \textbf{h}^{[1]}}}{\sum_{\text{words}} e^{W_{\text{word}} \cdot \textbf{h}^{[1]}}}
\ese

From here we just follow the backpropagation theory we have already developed, i.e.\ we calculate the update rules
using the chain rule, and we simply train the model through iteration of passing data. Notice however, that the
softmax function at the output is computationally very expensive since the denominator requires a summation over all
words in the vocabulary. We will revisit this issue soon.

\subsubsection{Continuous Skip-Gram}

\bd[Continuous Skip-Gram]
\textbf{Continuous Skip-Gram} is the reverse of CBOW, trying to predict the context of a word based on the target word.
\ed

\be
Considering our simple sentence from earlier, ``the man sat on a chair'', we get pairs of (context, word) where if we
consider a context window of size 1, we have examples like ([the], man), ([man], sat), ([sat], on) and so on. Now
considering that the skip-gram model's aim is to predict the context from the word, the model typically inverts the
contexts and words, and tries to predict each context word from its target word. Hence, the task becomes to predict
the context [the] given target word ``man'' and so on. Thus, the model tries to predict the context window words
based on the target word.
\vspace{-10pt}
\fig{skipgram}{0.4}
\ee

Notice that in the case of where the context window is equal to 1, the two models are in practice identical since
both of them try to predict one word based on another word. However, in the case where the window is greater than 1,
then CBOW tries to predict one word based on all the context words, while skip-gram tried to predict all the context
words based on one word. This is why we said that skip-gram is actually the reverse of CBOW\@. \v

Just like we discussed in the CBOW model, we need to model this skip-gram architecture now as a deep learning
classification model such that we take in the target word as our input and try to predict the context words. The
architecture is similar to the CBOW model with the difference that now everything is reversed. The input is again
just one word (e.g.\ the word ``on'') with one-hot representation:
\bse
\text{``on''} = [0,0,1,0, \ldots, 0]
\ese

but now the name of the weights are reversed since their meaning is reversed (see figure). So far only the naming has
changed. The mathematical difference is out the outcome, where now instead of just on word we want to predict all the
context words. We do that by having multiple predictions, hence multiple softmax output function and subsequently the
loss function is simply the sum of all individual cross-entropy loss functions:
\bse
\mathcal{L}(\theta) = - \sum_{\text{words in window}} \ln P(\text{word
\:} | \text{\: context})
\ese

From this point on everything is exactly the same. We train the model through backpropagation simply by computing the
update rules and running training iterations. Notice that the same problem as with CBOW is also present here, since
the softmax function at the output is computationally very expensive since the denominator requires a summation over
all words in the vocabulary. \v

There are 3 popular ways to overcome this problem. We will just mention them for now, but we will not explain them:
\bit
\item Solution 1: Negative sampling.
\item Solution 2: Contrastive estimation.
\item Solution 3: Hierarchical softmax.
\eit

Finally, it is worth mentioning the existence of a model that combines SVD and Word2vec called ``GloVe''.

\section{Positional Encoding}

\fig{trf03}{0.25}

The second step of the transformer architecture is the so called ``positional encoding''.

\bd[Positional Encoding]
\textbf{Positional encoding} is a simple embedding layer that takes as input the position of each element of the
input sequence, and it outputs a vector representation for each position.
\ed

The positional encoding layer is used to inject some information about the relative or absolute position of the
tokens in the sequence. The positional encoding vector is added to the embedding vector. Embeddings represent a token
in a d-dimensional space where tokens with similar meaning will be closer to each other. But the embeddings do not
encode the relative position of words in a sentence. So after adding the positional encoding, words will be closer to
each other based on the similarity of their meaning and their position in the sentence, in the d-dimensional space.
\v

\section{Attention}

Recall that we defined attention in the context of RNNs, as a mechanism that allows a model to focus on the relevant
parts of the input sequence at each time-step. The key difference between attention in RNNs and attention in
transformers is that in the latter, attention is not restricted to be between the hidden states of the encoder and
the decoder, but instead, it is expanded between any two elements of the input sequence. \v

First things first, let us begin by introducting some basic terminology that we will need in order to define attention
in a more abstract and generic way, to fit our needs for the context of transformers.

\bd[Query]
A \textbf{query} $Q$ is a quantity that we want to use to retrieve information from the input sequence.
\ed

\bd[Key]
A \textbf{key} $K$ is a quantity that we use to retrieve information from the input sequence.
\ed

\bd[Value]
A \textbf{value} $V$ is a quantity that we want to retrieve from the input sequence.
\ed

Based on the query and the key, we can define the concept of ``similarity'' between them, which correspons to the
concept of ``alignment score'' in the context of RNNs.

\bd[Similarity]
The \textbf{similarity} S between a query $Q$ and a key $K$ is a quantity that measures how similar the query is to
the key:
\bse
S(Q,K) = f(Q,K)
\ese
\ed

Once again, there are many functions $f$ that can be used to compute the similarity that have been proposed in the
literature. The most used one, and the one used in the paper of Vaswani et al.\ is the so-called ``scaled dot-product''
which is defined as:
\bse
S(Q,K) = \frac{QK^T}{\sqrt{d_k}}
\ese

where $d_k$ is the dimension of the query and the key. \v

Gienv the similarity we can now define the attention weights.

\bd[Attention Weights]
The \textbf{attention weights} $\alpha$ are the normalized similarities between a query $Q$ and a key $K$:
\bse
\alpha(Q,K) = \text{softmax}(S) = \text{softmax}\Big( \frac{QK^T}{\sqrt{d_k}} \Big)
\ese
\ed

Given the similarity, as we did in the context of RNNs, we can now define attention in the context of transformers.

\bd[Attention]
Given a query $Q$, a key $K$, and a value $V$, attention $\alpha$ is a vector defined as:
\bse
\boldsymbol{\alpha}(Q,K,V) = S \times V = \text{softmax}\Big( \frac{QK^T}{\sqrt{d_k}} \Big) \times V
\ese
\ed